% Generated draft: requires mathematical and source review.
\documentclass[11pt]{article}
\usepackage[margin=25mm]{geometry}
\usepackage{fontspec}
\setmainfont{Latin Modern Roman}
\usepackage{amsmath,amssymb,mathtools,unicode-math}
\setmathfont{Latin Modern Math}
\usepackage{enumitem,needspace,xurl,hyperref}
\setlength{\parindent}{0pt}
\setlength{\parskip}{5pt}
\setlength{\emergencystretch}{3em}
\newcommand{\N}{\mathbb{N}}
\newcommand{\Z}{\mathbb{Z}}
\newcommand{\Q}{\mathbb{Q}}
\newcommand{\R}{\mathbb{R}}
\newcommand{\C}{\mathbb{C}}
\newcommand{\F}{\mathbb{F}}
\newcommand{\A}{\mathbb{A}}
\newcommand{\PP}{\mathbb{P}}
\newcommand{\E}{\mathbb{E}}
\newcommand{\eps}{\varepsilon}
\newcommand{\dd}{\,\mathrm{d}}
\begin{document}
\section*{UnsolvedMath ID: 30004649}
\textbf{Source problem code:} catalogue entry
\subsection{Definitions and mathematical statement}
Let \(k/\Q\) be a finitely generated field extension, let \(\bar k\) be an algebraic closure, and put \(G_k=\operatorname{Gal}(\bar k/k)\). We use the conventional interpretation of a hyperbolic curve as a smooth, geometrically connected curve \(Y/k\) whose smooth proper compactification \(\overline Y\) satisfies
\[
2g(\overline Y)-2+\#(\overline Y_{\bar k}\setminus Y_{\bar k})>0.
\]
In particular, a proper hyperbolic curve has genus at least \(2\). These conventions make explicit hypotheses implicit in the supplied terminology.

For a geometric base point \(\bar y\), the étale fundamental group has an exact sequence
\[
1\longrightarrow\pi_1(Y_{\bar k},\bar y)\longrightarrow\pi_1(Y,\bar y)\xrightarrow{p_Y}G_k\longrightarrow1.
\]
A section is a continuous homomorphism \(s:G_k\to\pi_1(Y,\bar y)\) with \(p_Y\circ s=\operatorname{id}_{G_k}\). Write \(\operatorname{Sec}(Y/k)\) for sections modulo conjugation by \(\pi_1(Y_{\bar k},\bar y)\). A point \(y\in Y(k)\) induces a section class \([s_y]\), giving a map
\[
\kappa_Y:Y(k)\longrightarrow\operatorname{Sec}(Y/k).
\]
The use of conjugacy classes is necessary to make the point-induced section independent of choices of paths and base points; the supplied record does not explicitly specify this convention.

\textbf{Proper-curve statement (section conjecture).} For every finitely generated extension \(k/\Q\) and every proper hyperbolic curve \(X/k\), is \(\kappa_X\) bijective? Equivalently, does every continuous section \(s\) admit a unique \(x\in X(k)\) such that \([s]=\kappa_X(x)\)?

\textbf{Nonproper statement and unresolved qualification.} The source also asks whether, for every smooth nonproper hyperbolic curve \(U/k\) with compactification \(X\), every section of \(p_U\) is induced by a unique point of \(X(k)\). Here we interpret \(X\) as the smooth proper compactification and write \(j:U\hookrightarrow X\). For \(x\in U(k)\), the point induces a section class of \(U\) as above. For \(x\in X(k)\setminus U(k)\), there is no point morphism \(\operatorname{Spec}k\to U\), so the phrase ``induced by'' requires an additional definition at the boundary.

One precise possible interpretation, using the induced map \(j_*:\operatorname{Sec}(U/k)\to\operatorname{Sec}(X/k)\), is
\[
\forall\sigma\in\operatorname{Sec}(U/k),\quad
\exists!\,x\in X(k)\quad j_*(\sigma)=\kappa_X(x).
\]
This is explicitly a possible interpretation, not an asserted equivalent reformulation of the supplied nonproper claim. It only identifies the image section on the compactification. If the source intends a stronger assertion involving sections associated with boundary decomposition or inertia groups, the required definition and conditions are not supplied. Also, the compactification of a nonproper hyperbolic curve need not itself be hyperbolic, so the proper-curve conjecture above does not by itself cover every such \(X\). The intended meaning of the nonproper assertion therefore remains unresolved.
\subsection{Short English statement}
The proper section conjecture asks whether rational points of every proper hyperbolic curve over a finitely generated extension of \(\Q\) correspond bijectively to continuous Galois sections modulo geometric conjugacy. The source also includes nonproper curves and rational points of their compactifications, but does not define how a boundary point induces a section. A possible formulation using the compactification map is displayed with this limitation explicitly retained.
\subsection{Sources}
\begin{itemize}
\item Local UnsolvedMath registry, permanent ID \(30004649\), title: \emph{Sections and Rational Points of Hyperbolic Curves}; exact source problem code: \texttt{catalogue entry}. The record belongs to the Oberwolfach Reports Open Problems collection and supplies proposed year \(2021\), with no named proposer.
\item Supplied report attribution: \emph{Homotopic and Geometric Galois Theory (2022)}, Oberwolfach Reports, pp.\ 663--744, DOI \texttt{10.4171/owr/2021/12}: \url{https://doi.org/10.4171/owr/2021/12}. The literature list labels this same reference \emph{Oberwolfach Report 12/2021}. The differing year descriptions are retained without reconciliation.
\item Supplied literature source: \emph{On the birational section conjecture over finitely generated fields}, \url{https://doi.org/10.2140/ant.2021.15.435}.
\item Supplied literature assessment, checked on \texttt{2026-08-22}: the registry labels the problem partially solved and reports \(p\)-adic, birational, conditional, and special-case results, while describing the stated proper hyperbolic curve conjecture over finitely generated fields of characteristic zero as open. This is supplied dated curation derived from the OWR material, not an independently verified literature assessment. The record explicitly requests independent verification of time-sensitive status and cited claims.
\item The registry supplies the statement-status label \texttt{corrected\_verified} and reports a manual semantic, typographic, and LaTeX review. These are source metadata, not independent mathematical verification of this draft or resolution of its formulation ambiguities.
\end{itemize}
\end{document}
